\scanpage{154}
% archive foliation: 153.
\deleted{
\diagram{
\begin{tikzcd}[ampersand replacement=\&,column sep=large]
\&\& \unsure{G_0} \&\\
1 \arrow[r] \& N \arrow[r] \& G \arrow[r] \&
\unsure{\mathfrak S_A|\mathcal S_A} \arrow[ul,bend right=20] \arrow[r,dash] \arrow[d,phantom,"\wr"] \&1\\
\&\&\& \unsure{\mathfrak S_A|\mathcal S_A}
\end{tikzcd}
}
\[
\begin{gathered}
\unsure{\wr|N}\\
((\mathbb F_2)^A)'\\
\shortparallel\\
\operatorname{Ker}(\mathbb F_2^A-\mathbb F_2)\\
(x_i)\longmapsto\sum x_i
\end{gathered}
\]
\[
ux=yu\qquad
\begin{aligned}
y&=uxu^{-1}\\
&=u(x)
\end{aligned}
\]
\[
\begin{aligned}
g&=\underset{\underset{G_0}{\uparrow}}u.\underset{\underset N{\uparrow}}x=yu\\
g'&=u'.x'=y'u'
\end{aligned}
\qquad
g'g=y'\underbrace{u'yu'^{-1}}\underbrace{u'u}
\]
\[
\begin{aligned}
g'g&=u'x'ux
=(u'u)\underbrace{(u^{-1}x'u)x}_{u^{-1}(x')x}
=\underset{\underset{u'u}{\shortparallel}}{u''}
\underset{\underset{u^{-1}(x')x}{\shortparallel}}{x''}\\
&=(y'\unsure{u'|u}(y))u'u
\end{aligned}
\]
}
\hrule

L'extension quadratique de $\mathbb Q(t)$ est induite par
l'extension quadratique
\[
B=\mathbb Z[t]/(S^2-\unsure{P(t)|Q(t)})
\]
de $A=\mathbb Z[T]$. On voit que pour $\forall p\in\mathbb P\setminus\{2\}$,
localisant $A$ en $pA$ d'où $A_{pA}=\mathbb Z_{p\mathbb Z}\{t\}\overset{\mathrm{df}}=A_{(p)}$,
\deleted{\unsure{$B_{(p)}$}} $B_{(p)}\overset{\mathrm{df}}=B\otimes_A A_{(p)}$ est normal (que $p$ divise $q$, ou non).
Si $p=2$, alors $B_{(2)}$ est normal si $2\mid q$. Si $2\nmid q$, alors
il est normal aussi si $Q\bmod 2\in\mathbb F_2[t]$ n'est pas un carré
dans $\mathbb F_2[t]$ (donc pas non plus dans $\mathbb F_2(t)$). \deleted{\unsure{S'il}} Il en est
de $\unsure{=}$ s'il est un carré dans $\mathbb F_2[t]$, mais pas dans
$\mathbb Z/4\mathbb Z[t]$ --- i.e un certain invariant dans $\mathbb Z/2\mathbb Z[t]$ n'est pas nul.
Donc le seul cas où $B_{(2)}$ ne soit normal est celui où $2\nmid q$
\deleted{\unsure{et}} $Q(t)$ est un carré mod $4$ (p.ex.
\[
\begin{gathered}
Q(t)=t^2-4=(t+2)(t-2)\\
\text{ou}\quad Q(t)=4t+1,\quad\varepsilon=+1\\
\unsure{\text{Oui}|Q(t)}\ \unsure{\ne|+}\ 4t-1,\quad\varepsilon=-1
\end{gathered}
\]
)
Si on excepte ce cas, on voit que $B$ est normal en
codim $1$ --- i.e \inserted{$\operatorname{Spec}B$} n'a qu'un nb fini
\unsure{d'|de}\deleted{idéaux premiers} \inserted{pts non}
normaux, et ces pts sont fermés (donc isolés). Donc \deleted{\unsure{ce}} sauf
\deleted{l'est} au dessus d'un sous-schéma fini de $\operatorname{Spec}A=(\operatorname{Spec}\mathbb Z)[t]$,
$B$ représente la normalisée de $\mathbb E^1_{\mathbb Z}$ dans l'extension
quadratique envisagée --- et ceci vaut encore le cas
\unsure{d'une multiplicité}, si on remplace $\mathbb Z$ par $\mathbb Z[1/h]$.
